\begin{table}[htbp]
\centering
\caption{Required effective sample sizes under the rank-based model ($\rho \geq \targetRho$) and the localised Eppstein--Wang bound at different additive error tolerances ($\delta = 0.1$). The final column shows the normalised additive error that the rank model's sample size would guarantee under the EW bound.}
\label{tab:ew_comparison}
\small
\begin{tabular}{@{}rrrrrr@{}}
\toprule
\textbf{Reach} & \textbf{Rank $\effn$} & \textbf{EW $\varepsilon{=}0.05$} & \textbf{EW $\varepsilon{=}0.1$} & \textbf{EW $\varepsilon{=}0.2$} & \textbf{Implied $\varepsilon$} \\
\midrule
100 & 300 & 1,520 & 380 & 95 & 0.113 \\
500 & 300 & 1,842 & 461 & 115 & 0.124 \\
1,000 & 321 & 1,981 & 495 & 124 & 0.124 \\
3,000 & 556 & 2,200 & 550 & 138 & 0.099 \\
10,000 & 1016 & 2,441 & 610 & 153 & 0.078 \\
30,000 & 1760 & 2,661 & 665 & 166 & 0.061 \\
\bottomrule
\end{tabular}

\vspace{0.5em}
\footnotesize
The rank model requires substantially fewer samples than the EW bound at tight tolerances ($\varepsilon \leq 0.05$), but the two converge at $\varepsilon \approx 0.1$ for large reach. At reach above 3{,}000, the rank model implicitly guarantees $\varepsilon < 0.1$.
\end{table}
